\documentclass[11pt]{article}
\usepackage[utf8]{inputenc}
\usepackage[T1]{fontenc}
\usepackage[margin=1in]{geometry}
\usepackage{amsmath,amssymb}
\usepackage[colorlinks=true,urlcolor=blue,linkcolor=black,citecolor=black]{hyperref}

% Set the bibliography in a smaller font with tight inter-entry spacing (layout only; no content change).
\makeatletter
\let\oldthebibliography\thebibliography
\renewcommand{\thebibliography}[1]{%
  \small\oldthebibliography{#1}%
  \setlength{\itemsep}{0pt}\setlength{\parsep}{0pt}}
\makeatother

\title{A deterministic, auditable engine for cross-theory claim comparison:\\
cross-field discovery with built-in verification\thanks{Numbers locked to registry snapshot
\texttt{e9c877812c8c6897\ldots} (2026-10), pipeline v2 / corpus v1. Every figure renders from that one
snapshot; the locked registry and verification scripts accompany the public code release at
\url{https://github.com/Anuja-Khatri/consilience-verification}.}}

\author{Anuja Khatri\\[2pt]\small Independent Researcher}
\date{2026-10}

\begin{document}
\maketitle

\begin{abstract}
I present a deterministic engine that compares scientific claims across competing theories. Each claim is
first reduced to a fixed schema of typed fields, in the style of information extraction. The fields are its
entity type, its axis (the aspect of the system it addresses), its polarity (assertion versus negation),
and a value. A deterministic rule-based verifier, which I call the judge, then compares two claims from
these fields
alone. Its output is a verdict: either a grade of correspondence, which says how strongly the two claims
agree, or a fork, a question on which the theories make incompatible commitments. The judge is
reproducible, provable and auditable, and I re-express it three ways that agree with it exactly. As a logic
program, verdict equivalence is \textbf{1.0} ($n = 213$ pairs). As a first-order theory under an SMT
solver, congruence is \textbf{1.0} ($n = 311$ verdicts). As machine-checked proof objects, coverage is
\textbf{0.9842} ($n = 316$), with \textbf{0} integrity violations ($n = 316$). On the discovery side, the
engine recovers the COGITATE adversarial-collaboration key at \textbf{2 of 3} predictions (revised;
$n = 3$; not significant against a 1000-key permutation null, $p = 0.075$). The eLife synergistic-workspace
key recovery is below chance: \textbf{0 exact / 6 partial of 14} validated rows, with empirical
$p = 0.981$ ($n = 1000$; random keys recover more). The method's limits are first-class results. Both key
recoveries carry their chance rates, single-pass extraction carries a measured noise floor (its test-retest
reliability across repeated runs), and the extraction front-end under-propagates negation.
\end{abstract}

\section{Introduction}
Comparing scientific claims across competing theories is usually done by hand and by argument. This paper
presents a deterministic engine that does it mechanically and auditably. Each claim is reduced to a small
schema of typed fields (entity type, axis, polarity, value), in the style of information extraction. A
deterministic rule-based verifier, the judge, then compares two claims from those fields alone and returns
a grade of correspondence or a fork. Because the judge is a function of the typed fields, every comparison
is reproducible. The judge can also be re-expressed as a logic program, as a first-order theory under an
SMT solver \cite{z3}, and as machine-checked proof objects. This split, a language model proposing typed
candidates and a sound external program disposing, follows the LLM-Modulo pattern of a generator paired
with verifiers it cannot talk past \cite{kambhampati2024}. The judge reads only the typed fields, so the
generator has no text channel through which to argue a verdict into being.

I run the engine on a versioned, immutable dataset of six theories of consciousness (the frozen corpus).
Bridge extensions link active inference \cite{buckley2017} to machine learning: control as probabilistic
inference \cite{levine2018}, and predictive coding as approximate backpropagation \cite{millidge2022}. They
also span integrated information theory in its current formulation \cite{albantakis2023}. Nothing in the
method is specific to that domain. Throughout, the extraction front-end (its model, prompts and decoding
settings) is held fixed and versioned. I call this frozen pipeline the instrument.

I report two kinds of result. On the discovery side, the engine re-adjudicates a founding cross-field
correspondence record. It also tries to recover two published answer keys: the COGITATE adversarial
collaboration \cite{cogitate2025} and the eLife synergistic-workspace bridge \cite{luppi2024}. These keys
are held out from the engine's prompts, though they may sit in the models' pretraining. Each recovery is
scored against an explicit permutation chance rate. On the verification side, the judge's decisions are a
transparent, checkable function of the typed fields. A preregistered perturbation suite then measures how
the instrument behaves under meaning-preserving and meaning-flipping edits.

The verification lane is the contribution I stand behind. The discovery results are deliberately modest. I
report the method's limits as first-class results. Both key recoveries carry their chance rates, including
an honest below-chance negative. Single-pass extraction carries a measured noise floor. And the
under-propagation of negation is stated plainly. Every figure renders from a single numbers-locked registry
snapshot, so the same quantity never drifts from one document to the next.

\section{Discovery}

\subsection{Cross-field re-adjudication}
Before this evaluation, the engine's earliest corpus was a set of correspondences asserted across physics,
neuroscience and contemplative texts. I chose it deliberately. The vocabularies there are the most alien to
one another. So it is the hardest place for surface wording to do the work, and where the original claims
were made with the least discipline. The test is simple: if those early correspondences are re-typed
under today's frozen instrument (v2) and judged by today's fixed rules, how many survive? I re-ran every
pair end to end through the judge. Of the pairs with a recorded prior grade ($n = 7$), the new grade agrees
with the old one for a single pair (agreement \textbf{0.1429}, 1 of 7). It downgrades or withdraws the
other six (\textbf{0.8571}, 6 of 7). I read this as the method working, not failing. Most early cross-field
correspondences do not survive re-typing under a fixed codebook (the judge's rule set and field
definitions). A judge strict enough to retire most of its builder's own early claims is one whose surviving
claims are worth attention.

\subsection{Unification index}
The unification index is a ranking. It orders the correspondences the engine has kept (the claim pairs the
judge graded as matching) by how much unifying work each one does. The pairs that tie the most theories
together rise to the top. A ranking like this is useful only if it is stable. The worry is that the order
is an artifact of the weights I chose. So the weight-simplex test recomputes the ranking under every
weighting across the simplex and asks whether the order holds. I scored each correspondence by a
preregistered mix of three ingredients, with weights 1/1/1. The ingredients are how many forks it touches,
the size of its supporting evidence cone, and how many axes it reaches. I then swept the weights and measured Kendall's
$\tau$ \cite{kendall1938}, the rank correlation between the base ranking and each reweighted one. Across the
216 correspondences, the smallest $\tau$ over the sweep is \textbf{0.7165} ($n = 216$). So the order is
moderately weight-sensitive, not robust: largely but not entirely weight-independent, and one ingredient
dominates the score, which I flag.

\subsection{Held-out key recoveries, with chance rates}
COGITATE is an adversarial collaboration between proponents of two rival theories of consciousness:
integrated information theory and global neuronal workspace theory (IIT and GNWT, respectively). Over
several years the two camps negotiated a shared set of predictions in advance, together with the
experiments that would decide between them \cite{cogitate2025}. That key is the right benchmark precisely
because it was so expensive to produce. Suppose the engine reads only the theories' own texts and still
recovers predictions that took years of expert argument to fix. Then the judge is doing something the
surface text alone does not give it. The key was preregistered before the collaboration's results were known
\cite{cogitate_prereg}, and the engine never saw the key text, so the holdout is clean. I ran the
instrument end to end against the three held-out predictions, then applied the review gate, which accepts
or rejects each extracted claim. The engine recovers \textbf{2 of 3} predictions; the full counts of
recovered, partial, missed and beyond-key claims are \{2,0,1,9\} ($n = 3$). Outcome-informed disclosure
(verbatim): ``iit\_c009 was initially rejected at the gate; after a sensitivity check showed it was the
sole binding claim for P3, it was re-reviewed and approved on codebook-consistency grounds (matching
approved canonical iit\_c025). Original gated result \{1,0,2,9\} and revised \{2,0,1,9\} are both reported;
fully ungated sensitivity \{2,0,1,16\}.'' The one remaining miss is the second prediction (P2, on temporal
profile), and it is a property of the instrument, not the gate: no IIT temporal\_profile claim was ever
extracted. Is recovering two of three meaningful? I compared it against a null of 1000 random typed-row
keys. The empirical permutation $p$ is the share of those keys that recover at least as much as the real
one. Here it is \textbf{$p = 0.075$} ($n = 1000$), beside a normal approximation $z = 2.2693$ from the same
null.
But that null is discrete and right-skewed, so the $z$ overstates significance and the empirical
\textbf{$p = 0.075$} ($n = 1000$) governs. It does not reach $\alpha = 0.05$, so I read the COGITATE
recovery as suggestive only.

The eLife bridge is a second, independent key of a different kind. It is a published, information-theoretic
analysis that builds a quantitative bridge between the same two theories, deriving where integrated
information and the global workspace should line up \cite{luppi2024}. Different authors built it with
different machinery than COGITATE, so it is an independent check against a key no one on the engine side
wrote. I compared the engine's IIT$\times$GNWT fork and correspondence set against the \textbf{14 validated}
bridge rows. The engine recovers \textbf{0 exact} matches and \textbf{6 partial} ones ($n = 14$). It misses
8, and \textbf{0} rows are unmappable to any engine axis ($n = 14$). (The provisional build had 16 rows; 2
were dropped at the format gate as methods, including the one vocabulary-gap row. So the validated
denominator of 14 carries 0 unmappable.) Scored against the same 1000-key null, the result does not beat
chance. Random keys recover more on average (null mean 9.2). The share of null keys doing at least as
well as the engine is empirical \textbf{$p = 0.981$} ($n = 1000$), below the null. So I report it as found:
on the engine's IIT$\times$GNWT set, the eLife bridge's partial recovery is at or below chance.

\section{Verification}

\subsection{Audit method}
Which of the engine's own earlier discoveries survive a strict audit? Before trusting any correspondence
the engine surfaces, I wanted to know how many of its prior demonstration discoveries hold up under
explicit criteria rather than impression. I ran a five-part audit on each: provenance, quote-support, a
screen for superficial analogy (the Newman screen), testability, and scale. A discovery fails if it misses
even one. Of the 9 demonstration discoveries, \textbf{1 of 9} survives all five ($n = 9$). Of a separate
sample of 30 engine-asserted matches, \textbf{2 of 30} survive ($n = 30$). The audit is deliberately
strict, and most prior discoveries fail at least one criterion. I treat that as the audit doing its job.

\subsection{Instrument}
How much does the extraction front-end disagree with itself when the input has not changed? This is its
test-retest reliability, its self-consistency across repeated runs, which I report as the stability. Every
downstream invariance and tracking figure is only meaningful against that floor. So I fix one estimate of
it and use it throughout, scoring each contrast against the measured floor rather than against an
assumption of zero noise. The adopted decoding is \texttt{baseline-single-pass}. I estimate the stability
with the primary estimator over repeated runs of the same passages. Across the 88 repeated-run units, the
stability is \textbf{0.25} (95\% CI [0.1713, 0.3496], $n = 88$). For comparison, an earlier, coarser
estimator on the same 80 runs gives 0.271. A noise floor this high is itself a limitation, and every
invariance number below is read against it.

\subsection{The verified judge}
Can the judge's rules be proved correct, not just trusted? The judge is a function of the typed fields
alone. So independent formal re-expressions of the same rules should reproduce its verdicts exactly. That would
make each verdict a checkable consequence of the inputs, not the behaviour of one program. I
re-expressed the judge three ways and checked each against it over the corpus. The three forms are a logic
program, a first-order theory discharged by the Z3 SMT solver \cite{z3}, and machine-checked proof objects
in the spirit of formalised-mathematics libraries \cite{mathlib}. Across the 213 theory pairs, the logic
program
returns the identical verdict every time (verdict equivalence \textbf{1.0}, $n = 213$). Under the solver,
the first-order theory is congruent with the judge on every verdict checked (\textbf{1.0}, $n = 311$). The
proof objects cover \textbf{0.9842} of the 316 derivations ($n = 316$), with \textbf{0} integrity
violations across the same 316 ($n = 316$). Three independent re-expressions agreeing exactly is what the
verification lane rests on: the verdicts are provable from their inputs, not taken on trust.

\subsection{Statement-level contracts: the negation headline}
What happens to a verdict when a sentence's meaning is flipped by negation? The judge is only trustworthy
if the fields it reads track meaning. So I ran a preregistered, statement-level perturbation suite. It
applies meaning-preserving and meaning-flipping edits to each passage and measures how the typed fields and
the verdicts respond, scored against a re-extraction noise-floor control (OSF preregistration DOI
\href{https://doi.org/10.17605/OSF.IO/5N24M}{10.17605/OSF.IO/5N24M}). One caveat frames the polarity
numbers. The claim store, the versioned, immutable dataset of extracted claims, is nearly single-class on
polarity, that is on assertion versus negation (92 asserts, 3 denies). Because of that imbalance, the
polarity stability under an identity control (the same passage re-extracted with no edit) is \textbf{0.975}
($n = 40$). The panel polarity agreement below mostly reflects the extractor's default of assertion, not a
demonstrated reading of polarity. So the negation figure is the informative one. Under explicit negation,
the extracted polarity flips as it should only \textbf{0.1765} of the time ($n = 34$). Under double
negation, the extractor restores the assertion \textbf{0.0} of the time ($n = 8$, a low count because the
perturber rarely negates twice). And the verdict itself tracks the edit correctly for \textbf{0.4189} of
the 74 meaning-flipping edits ($n = 74$). The judge is correct given its fields. What these numbers locate
is the extraction front-end under-propagating negation into those fields. A single-reviewer hand-check
agrees with the judge's automated verification on \textbf{0.893} of 28 items ($n = 28$). So the gap is in
extraction, not in the check.

\subsection{Obfuscation control}
Does hiding each theory's coined vocabulary change the engine's verdicts? If the correspondences were
driven by surface wording, masking the signature terms should degrade them. If they are driven by the typed
fields, masking should leave them intact. I ran two arms of extraction and fork recovery, identical except
that one worked on passages with the theories' surface vocabulary masked. I recorded both the change in
recovery and whether the extractor re-emitted the coined terms. Fork recovery is the share of the known
forks the engine reproduces. Across the 14 reference forks, masking moves it by only \textbf{0.072}
($n = 14$). So surface masking does not reduce fork recovery. But it does not blind the model to theory
identity either. Over those same 14, the extractor re-identifies the theories and re-emits their coined
terms, producing \textbf{12} such leaks ($n = 14$). The verdicts survive masking, but the model still knows
which theories it is reading.

\subsection{Reliability framing}
Do two different vendors' models type the same claims the same way? Cross-vendor agreement tells us whether
the typing protocol is reproducible across independent models, not whether its verdicts are valid. I ran a
two-vendor panel (gpt-4.1 $\times$ gemini-3.1-pro-preview, full coverage) over the 95 claims typed by both.
The two models agree moderately on entity type (Cohen's $\kappa$ \cite{cohen1960} \textbf{0.525}, $n = 95$)
and on axis (Cohen's $\kappa$ \textbf{0.503}, $n = 95$). They agree well on polarity (PABAK, the
prevalence- and bias-adjusted kappa \cite{byrt1993}, \textbf{0.811}, $n = 95$), though that figure carries
the single-class caveat from the negation headline. Within a single model, the stability is \textbf{0.25}
(the primary estimator, 95\% CI [0.1713, 0.3496], $n = 88$; an earlier, coarser estimator gives 0.271). So
the protocol is reproducible across models. That is a statement about reproducibility, not about whether the
typing is valid.

\section{Limitations}
These limitations carry over from an earlier internal evaluation that this paper supersedes. The engine's
six theories are all theories of consciousness, so the results may not transfer to other fields. The
cross-field reach rests on the re-adjudication record and the corpus extensions, not the core. The
reference labels were made by the builder, so the judge-after-audit reference reflects one builder's
choices. A second, independent annotator is planned but not yet run, so human agreement, the ceiling
against which model agreement is read, is not yet available. In production, a smaller same-family model
cross-checks the main extractor; the two-vendor panel is an evaluation add-on that measures reproducibility,
not validity. The corpus is small (95 approved claims, six theories). The deeper arms, embedding alignment,
unsupervised optimal transport via Gromov-Wasserstein couplings \cite{memoli2011}, and partial-information
decomposition of the verdict \cite{williams_beer2010}, are pilot-scale feasibility studies in a planned
second paper.

Single-pass extraction carries a measured noise floor (stability \textbf{0.25}, $n = 88$), so every
invariance and tracking number is read relative to the identity control. Negation tracking is itself a
measured limitation. The store is nearly single-class on polarity (92 asserts / 3 denies), so the
polarity-stability figures reflect the extractor's default, not a demonstrated reading of polarity. Holdout
controls the prompt, not the training data: the held-out keys never entered the extraction or judging
prompts (measured clean), but they may sit in a model's pretraining. The obfuscation control is what
addresses that, and it shows identity is not blinded. Several review packets were judged at the statement
level, or from truncated excerpts without full fork-side source text, so those arms are reported as ``no
fork flagged by the scan'', not as validation of the forks. Decisions were model-drafted suggestions
(Claude) reviewed and confirmed by the author. And the open-weight panel leg is not run (no
hosted-open-weight key), so the panel is two-vendor.

\section{Conclusion}
The engine turns cross-theory claim comparison into a reproducible, auditable computation: a typed
intermediate representation, a deterministic judge, and three independent re-expressions of it that agree
exactly. Two lessons generalise beyond this corpus. First, a generator-and-verifier split pays off when the
verifier is provable and reads only the structured fields. Then the generating model has no text channel
through which to argue a verdict into being. Second, honest evaluation here rested on fixed reference
points. The perturbation suite was preregistered and scored against a measured noise floor, not against
zero. Chance rates came from random-key permutation, so that even a below-chance recovery is reported as
one. And every document renders from a single locked registry snapshot, so no number drifts from one place
to the next. The discovery claims are modest by construction: most prior cross-field correspondences do not
survive re-typing, one held-out key is recovered only suggestively, and the other below chance. I report
them that way. The verification lane, not the discovery headline, is the contribution I stand behind: a
claim-comparison verdict that is a checkable function of its inputs.

\bibliographystyle{plain}
\bibliography{refs}

\end{document}
